% MEC Lab 1. Compile: latexmk -xelatex lab01.tex
\documentclass[10pt,aspectratio=169,t]{beamer}
\usetheme[lab]{upbminimal}

\course{Mobile and Embedded Computing}
\decklabel{Lab 1}
\title{Toolchain, board, broker, and project kickoff}
\subtitle{TinyGo, a blinking board, a local broker, and your team}
\author{Dinu-Ștefan Rusu}
\institute{FILS, Universitatea Politehnica București}
\date{Week 2}

\begin{document}

\titleframe

\begin{frame}[eyebrow=Today]{What this lab is for}
  \vskip 2mm
  \begin{grid}[2]
    \cell[\faMicrochip]{Put code on a device}{Install TinyGo, build the blink program for your board or for the simulator, flash it, and read what it prints on the serial console.}
    \cell[\faServer]{Run your own broker}{Install Mosquitto, publish and subscribe from the command line, then agree on a project idea that meets all five requirements.}
  \end{grid}
  \vskip 8mm
  \taskmeta{Time}{2 hours}{Submit}{Branch \code{lab-1} in your team repository, before next lab}
  \vskip 3mm
  \muted{Bring a board if you have one. If you do not, Task III replaces Task II and is worth the same.}
\end{frame}

\begin{frame}[fragile,eyebrow=Setup]{Where your work goes}
  \begin{steps}
    \item One repository per team. Use the Flutter app repository from ADMD: the project extends that app.
    \item Create the branch and the folder this lab writes into:
\begin{shell}
git checkout -b lab-1
mkdir -p device/blink
\end{shell}
    \item Open a pull request from \code{lab-1} on the first commit. Everything you are graded on goes in its description.
  \end{steps}
  \begin{hint}
    Pull requests, not pushes to main. Your individual mark is read from your own commits.
  \end{hint}
\end{frame}

\begin{frame}[fragile,eyebrow=Setup]{Install Go and TinyGo}
  \vskip 1mm
\begin{shell}
# macOS
brew install go
brew tap tinygo-org/tools && brew install tinygo

# Windows, with Scoop
scoop install go tinygo

# Debian and Ubuntu: Go from go.dev/dl, TinyGo from its releases page
sudo dpkg -i tinygo_*.deb
tinygo version
\end{shell}
  \muted{Install Go first: TinyGo reuses the Go standard library and the module cache. If \code{tinygo version} is not found, the TinyGo \code{bin} folder is not on your PATH.}
\end{frame}

\begin{frame}[fragile,embedded,eyebrow=Task I]{Pick a board and confirm its target}
  \begin{columns}[T]
    \begin{column}{0.55\textwidth}
      \begin{steps}
        \item Pick one: a Raspberry Pi Pico, an ESP32-C3 board, or an ESP32-S3 board. With no hardware, go to Task III.
        \item List what your TinyGo can build for:
\begin{shell}
tinygo targets
\end{shell}
        \item Take the target name from your board's page on the TinyGo site. The Pico is \code{pico}, a plain ESP32-C3 board is \code{esp32c3}.
      \end{steps}
    \end{column}
    \begin{column}{0.41\textwidth}
      \kv{Done when}{}
      \begin{checklist}
        \item \code{tinygo version} prints a version.
        \item Your target appears in the output of \code{tinygo targets}.
        \item The target name is written in the pull request description.
      \end{checklist}
      \begin{warning}
        Board variants have their own targets, such as \code{xiao-esp32c3}. Check the board page for the exact name.
      \end{warning}
    \end{column}
  \end{columns}
\end{frame}

\begin{frame}[fragile,embedded,eyebrow=Reference]{The blink program, in full}
\begin{golang}
package main
import "machine"
import "time"
func main() {
    led := machine.LED
    led.Configure(machine.PinConfig{Mode: machine.PinOutput})
    for {
        led.Set(!led.Get())
        println("blink")
        time.Sleep(500 * time.Millisecond)
    }
}
\end{golang}
  \muted{Save it as \code{device/blink/main.go}. Use the LED pin your board page names.}
\end{frame}

\begin{frame}[fragile,embedded,eyebrow=Task II]{Flash the board and read the console}
  \begin{columns}[T]
    \begin{column}{0.55\textwidth}
      \begin{steps}
        \item Plug the board in with a data cable, not a charge-only cable.
        \item Build and flash in one command:
\begin{shell}
tinygo flash -target=pico ./device/blink
\end{shell}
        \item Watch the LED, then read what the program prints:
\begin{shell}
tinygo monitor
\end{shell}
        \item Stop the monitor with Ctrl-C.
      \end{steps}
    \end{column}
    \begin{column}{0.41\textwidth}
      \kv{Done when}{}
      \begin{checklist}
        \item The LED blinks about once a second.
        \item \code{tinygo monitor} prints one line per blink.
        \item A photo or a short video of the board is in the pull request.
      \end{checklist}
    \end{column}
  \end{columns}
\end{frame}

\begin{frame}[fragile,embedded,eyebrow=Task III]{No board: the same program in Wokwi}
  \begin{columns}[T]
    \begin{column}{0.55\textwidth}
      \begin{steps}
        \item On wokwi.com open the ESP32 custom application template, then set the board to ESP32-C3 DevKitM-1.
        \item Add an LED and a 220 ohm resistor from the parts menu, wired to a free GPIO pin. Use that pin in the code.
        \item Build the firmware on your laptop:
\begin{shell}
tinygo build -o blink.elf \
  -target=esp32c3 ./device/blink
\end{shell}
        \item Press F1, choose \emph{Upload Firmware and Start Simulation}, and pick \code{blink.elf}.
      \end{steps}
    \end{column}
    \begin{column}{0.41\textwidth}
      \kv{Done when}{}
      \begin{checklist}
        \item The simulated LED blinks.
        \item The Wokwi serial monitor prints one line per blink.
        \item One screenshot shows the board and the serial output.
      \end{checklist}
    \end{column}
  \end{columns}
\end{frame}

\begin{frame}[fragile,eyebrow=Setup]{Install the Mosquitto broker}
\begin{shell}
# macOS
brew install mosquitto

# Debian and Ubuntu
sudo apt install mosquitto mosquitto-clients

mosquitto -v         # foreground, port 1883, one line per event
\end{shell}
  \begin{hint}
    Mosquitto 2 and later accept local connections only until you configure a listener. In this lab that default is what you want.
  \end{hint}
\end{frame}

\begin{frame}[fragile,eyebrow=Task IV]{Publish and subscribe from two terminals}
  \begin{columns}[T]
    \begin{column}{0.55\textwidth}
      \begin{steps}
        \item Leave the broker running. In a second terminal, subscribe to a wildcard topic:
\begin{shell}
mosquitto_sub -h localhost -t mec/lab1/#
\end{shell}
        \item In a third terminal, publish to a subtopic:
\begin{shell}
mosquitto_pub -h localhost \
  -t mec/lab1/hello -m "1"
\end{shell}
        \item Publish three more messages, changing the subtopic each time.
      \end{steps}
    \end{column}
    \begin{column}{0.41\textwidth}
      \kv{Done when}{}
      \begin{checklist}
        \item The subscriber prints every message the publisher sends.
        \item Messages from different subtopics all reach the wildcard subscription.
        \item The broker log lists both clients connecting.
      \end{checklist}
    \end{column}
  \end{columns}
\end{frame}

\begin{frame}[eyebrow=Task V]{Your team and your project idea}
  \begin{columns}[T]
    \begin{column}{0.55\textwidth}
      \begin{steps}
        \item Form a team of at most three. Fill one row in the team sheet in the Teams channel: names, GitHub handles, repository URL.
        \item Write the idea in the repository README: one paragraph, then one line for each of the five requirements on the next slide.
        \item Show it to me during this lab. Milestone 1 is the team and the idea approved.
      \end{steps}
    \end{column}
    \begin{column}{0.41\textwidth}
      \kv{Done when}{}
      \begin{checklist}
        \item The sheet row is filled and the repository URL opens.
        \item The README names the sensor, the transport, and the conflict rule.
        \item The idea is approved.
      \end{checklist}
    \end{column}
  \end{columns}
\end{frame}

\begin{frame}[eyebrow=Reference]{The five requirements your idea must meet}
  {\usebeamerfont{small}
  \begin{tabular}{@{}p{0.30\textwidth}p{0.62\textwidth}@{}}
    \toprule
    \thead{Requirement} & \thead{True at the defense} \\
    \midrule
    Device component & A TinyGo program on a board or in Wokwi, reading a sensor. \\
    Device to phone & The reading arrives over MQTT or BLE and the app shows it live. \\
    Offline-first & The app works with no network, syncs later, states its conflict rule. \\
    One boundary & Either gRPC with a Go server, or FFI, or a platform channel. \\
    A measured claim & One number your team measured, with the method. \\
    \bottomrule
  \end{tabular}}
  \begin{takeaway}{All five, or the project is not approved.}
    Pick an idea where each requirement has an obvious home.
  \end{takeaway}
\end{frame}

\begin{frame}[embedded,eyebrow=Troubleshooting]{What goes wrong, and what to do}
  \vskip 2mm
  \trouble{permission denied: /dev/ttyUSB0}{Your user is not in the serial port group. Add yourself to \code{dialout} on Debian and Ubuntu, or \code{uucp} on Arch, then log out and back in. Do not flash with sudo.}
  \trouble{no serial ports found}{The cable carries power only, or the board is not in its bootloader. Swap the cable, then hold the BOOT button while plugging the board in and release it.}
  \trouble{unknown target or board name}{The string after \code{-target=} is not one TinyGo knows. Copy it from \code{tinygo targets} or from the board's page.}
\end{frame}

\begin{frame}[eyebrow=Troubleshooting]{When the broker refuses you}
  \vskip 2mm
  \trouble{Error: Connection refused}{Nothing is listening on port 1883. Start the broker in a terminal with \code{mosquitto -v} and leave it running while you test.}
  \trouble{Connection refused from another machine}{The default configuration binds to localhost. Add \code{listener 1883} and \code{allow\_anonymous true} to \code{mosquitto.conf}, and open the port in the local firewall. Lab network only.}
  \trouble{Address already in use}{A broker is already running, often as a system service. Use that one, or stop the service before starting your own.}
\end{frame}

\begin{frame}[eyebrow=Submission]{What the pull request has to contain}
  \vskip 2mm
  \begin{grid}[2]
    \cell[\faMicrochip]{Proof it ran}{A photo of the blinking board, or a screenshot of the Wokwi simulation, plus the serial output copied as text.}
    \cell[\faServer]{Proof it talks}{The two terminals from Task IV, subscriber and publisher, with the messages visible.}
  \end{grid}
  \vskip 4mm
  \kv{Also in the description}{The board and the target name you used, and a link to the team sheet row.}
  \begin{takeaway}{Open the pull request now, not at the deadline.}
    A branch with no pull request is not a submission, and an empty description is not evidence.
  \end{takeaway}
\end{frame}

\closingframe{Push before you leave.}{Branch lab-1 · blink and serial output in the PR · idea approved}

\end{document}
